\subsection{The Grothendieck Group R(A)}

Let A be a ring. Let \(\mathcal{F}(A)\) be the set of classes of finitely generated A-modules (VIII, p.~51). The classes of the A-modules of finite length form a subset \(\mathcal{L}\mathcal{F}(A)\) of \(\mathcal{F}(A)\) ; we have seen that \(\mathcal{L}\mathcal{F}(A)\) is a hereditary set of classes of modules. We denote the Grothendieck group associated with \(\mathcal{L}\mathcal{F}(A)\) by R(A) and the image in R(A) of the class of an A-module E of finite length by {[}E{]}.

The results of Subsection 3 imply the following:

\begin{enumerate}
\def\labelenumi{\alph{enumi})}
\item
  Let \(\mathcal{S}(\mathrm{A})\) be the set of classes of simple A-modules. The family \(([S])_{S\in \mathcal{S}(A)}\) is a basis of the \(\mathbf{Z}\) -module \(R(A)\) .
\item
  Let \(\mathrm{E}\) and \(\mathrm{F}\) be semisimple A-modules of finite length. Then \(\mathrm{E}\) and \(\mathrm{F}\) are isomorphic if and only if we have \([\mathrm{E}] = [\mathrm{F}]\) in \(\mathrm{R(A)}\) .
\item
  Let E be an A-module of finite length and \((\mathrm{E}_{i})_{0\leqslant i\leqslant n}\) a Jordan--Hölder series of E. Set \(F=\bigoplus_{i=1}^{n}(\mathrm{E}_{i-1}/\mathrm{E}_{i})\) . Then F is a semisimple A-module of finite length, and we have \([E]=[F]\) in R(A).
\item
  Let \(\ell : \mathrm{R}(\mathrm{A}) \to \mathbf{Z}\) be the homomorphism characterized by \(\ell([\mathrm{S}]) = 1\) for every simple A-module S. We then have \(\ell([\mathrm{E}]) = \sum_{\mathrm{S} \in \mathcal{S}(\mathrm{A})} \ell_{\mathrm{S}}(\mathrm{E}) = \mathrm{long}_{\mathrm{A}}(\mathrm{E})\) for every A-module E of finite length.
\end{enumerate}

If D is a field, then the homomorphism \(\ell: \mathrm{R}(\mathrm{D}) \to \mathbf{Z}\) is an isomorphism.

Remarks. --- 1) Let \(\mathcal{S}\mathcal{S}(\mathrm{A})\) be the hereditary set of classes of semisimple A-modules of finite length. By Proposition 7 of VIII, p.~190, the Grothendieck group \(\mathrm{K}(\mathcal{S}\mathcal{S}(\mathrm{A}))\) is a free \(\mathbf{Z}\) -module whose elements {[}S{]} \(_{\mathcal{S}\mathcal{S}(\mathrm{A})}\) (for \(\mathrm{S} \in \mathcal{S}(\mathrm{A})\) ) form a basis. The canonical homomorphism \(\gamma_{\mathcal{L}\mathcal{F}(\mathrm{A}),\mathcal{S}\mathcal{S}(\mathrm{A})}\) (VIII, p.~186) is therefore an isomorphism.

\begin{enumerate}
\def\labelenumi{\arabic{enumi})}
\setcounter{enumi}{1}
\item
  Let \(\mathrm{K}^{\prime}(\mathcal{L}\mathcal{F}(\mathrm{A}))\) be the Grothendieck group of \(\mathcal{L}\mathcal{F}(\mathrm{A})\) for direct sums (VIII, p.~188). The Krull--Remak--Schmidt theorem (VIII, p.~37) implies that \(\mathrm{K}^{\prime}(\mathcal{L}\mathcal{F}(\mathrm{A}))\) is a free Z-module with basis the set of classes of indecomposable A-modules of finite length, while \(\mathrm{K}(\mathcal{L}\mathcal{F}(\mathrm{A}))\) admits the set of classes of simple A-modules as a basis.
\item
  Let E be an A-module of finite length. By c), there exists a semisimple A-module \(E'\) of finite length such that \([E] = [E']\) , and by b), such a module is defined up to an isomorphism; we sometimes call it a semisimplification of E.
\end{enumerate}

PROPOSITION 8. --- Let A be a principal ideal domain that is not a field, and let L be its field of fractions. There exists an isomorphism \(\varphi : \mathrm{R}(\mathrm{A}) \to \mathrm{L}^{*}/\mathrm{A}^{*}\)

such that we have

\[
\varphi ([ \mathrm{A} / a \mathrm{A} ]) = a \mathrm{A} ^ {*}\tag{9}
\]

for every \(a \neq 0\) in A.

Let P be a system of representatives consisting of irreducible elements of A (VII, §1, No.~3, p.~3). The maximal ideals of A are the ideals pA for \(p \in P\) ; every simple A-module is therefore isomorphic to a unique module A/pA. Moreover (VII, §1, No.~3, p.~4, Theorem 2), the abelian group \(L^{*}/A^{*}\) is free and has the family \((pA^{*})_{p \in P}\) as a basis. Hence there exists an isomorphism \(\varphi\) from R(A) to \(L^{*}/A^{*}\) characterized by \(\varphi([A/pA]) = pA^{*}\) for every \(p \in P\) .

Let \(a \neq 0\) in A. There exist an integer \(r \geqslant 0\) , elements \(p_1, \ldots, p_r\) of P, and an element \(u\) of A* such that we have \(a = up_1 \cdots p_r\) . The module A/aA admits the composition series defined by

\[
\mathrm{E} _ {0} = \mathrm{A} / a \mathrm{A}, \quad \mathrm{E} _ {i} = (p _ {1} \dots p _ {i} \mathrm{A}) / a \mathrm{A} \quad (1 \leqslant i \leqslant r),
\]

and the module \(\mathrm{E}_{i-1}/\mathrm{E}_{i}=(p_{1}\cdots p_{i-1}\mathrm{A})/(p_{1}\cdots p_{i}\mathrm{A})\) is isomorphic to \(A/p_{i}A\) . We therefore have (VIII, p.~185, Proposition 3)

\[
\varphi ([ \mathrm{A} / a \mathrm{A} ]) = \prod_ {i = 1} ^ {r} \varphi ([ \mathrm{A} / p _ {i} \mathrm{A} ]) = p _ {1} \dots p _ {r} \mathrm{A} ^ {*} = a \mathrm{A} ^ {*}.
\]

Remarks. --- 4) We keep the assumptions and notation of Proposition 8. Let E be an A-module of finite length, and let \(a_{1}\mathrm{A},\ldots ,a_{n}\mathrm{A}\) be its invariant factors (VII, §4, No.~5, p.~20). Since E is isomorphic to \(\bigoplus_{i = 1}^{n}\mathrm{A} / a_i\mathrm{A}\) , we have

\[
\varphi ([ \mathrm{E} ]) = \prod_ {i = 1} ^ {n} \varphi ([ \mathrm{A} / a _ {i} \mathrm{A} ]) = a _ {1} \dots a _ {n} \mathrm{A} ^ {*}.\tag{10}
\]

\begin{enumerate}
\def\labelenumi{\arabic{enumi})}
\setcounter{enumi}{4}
\tightlist
\item
  *Let A be a Dedekind ring that is not a field (Comm. Alg., VIII, §2, No.~1). By reasoning as in Proposition 8, we can prove the existence of an isomorphism \(\varphi\) from R(A) to the group of fractional ideals of A, characterized by \(\varphi([A/\mathfrak{a}]) = \mathfrak{a}\) for every nonzero ideal \(\mathfrak{a}\) of A.*
\end{enumerate}

Proposition 8 will be used, for example, in the following two cases:

\begin{enumerate}
\def\labelenumi{\alph{enumi})}
\tightlist
\item
  Suppose that we have A = Z. The Z-modules of finite length are simply the finite abelian groups. Since \(Q^{*}\) is the direct product of \(Z^{*} = \{1, -1\}\) and \(Q_{+}^{*}\) , we can deduce from Proposition 8 an isomorphism \(\varphi'\) from \(\mathrm{R}(\mathbf{Z})\) to \(Q_{+}^{*}\) given by
\end{enumerate}

\[
\varphi^ {\prime} ([ \mathrm{G} ]) = \operatorname{Card} (\mathrm{G})
\]

for every finite abelian group G.

\begin{enumerate}
\def\labelenumi{\alph{enumi})}
\setcounter{enumi}{1}
\tightlist
\item
  Suppose that A is the polynomial ring K{[}T{]} in one variable T with coefficients in a commutative field K. Let E be a finite-dimensional vector space over K and u an endomorphism of E. As in VII, §5, No.~1, p.~29, denote by \(E_{u}\) the A-module with underlying additive group E and external law \((p, x) \mapsto p(u)x\) . The A-module \(E_{u}\) has finite length. Conversely, every simple A-module is finite-dimensional over K (VII, §4, No.~8, p.~25, Remark 4). Consequently, every A-module of finite length is finite-dimensional over K, hence of the form \(E_{u}\) . Moreover (VII, §5, No.~3, p.~32, Corollary 1), the product of the invariant factors of \(E_{u}\) is equal to the characteristic polynomial \(\chi_{u}\) of u. Consequently, Proposition 8 gives an isomorphism
\end{enumerate}

\[
\varphi : \mathrm{R} (\mathrm{K} [ \mathrm{T} ]) \to \mathrm{K} (\mathrm{T}) ^ {*} / \mathrm{K} ^ {*}
\]

characterized by \(\varphi([\mathrm{E}_u]) = \chi_u\mathrm{K}^*\) (cf.~formula (10)).

